\documentclass{article}
\usepackage{tikz}
\usetikzlibrary{angles, quotes}
\usepackage{amssymb}
\usepackage{pgfplots}
\pgfplotsset{width=10cm,compat=1.9}
\usepgfplotslibrary{statistics}
\usepackage[left=2.54cm,right=2.54cm,top=2.54cm,bottom=2.54cm]{geometry}
\usepackage{tasks} %Paquete necesario para la producción de items horizontales para algunos ejercicios de tarea empleando el comando /begin{multicols}}{#}.../end{multicols} para ayudarnos a reducir las páginas
\usepackage{pdflscape} %comando que permite cambiar a orientación horizontal las hojas%
\usepackage{soul} %Paquete para permitir la compilación de acentos usando el comando {\hl{}}}\\
\sethlcolor{yellow(munsell)} %comando para definir el color para el subrayado usando el comando \hl
\usepackage[utf8]{inputenc}
\usepackage[spanish]{babel}
\usepackage{icomma}
\usepackage{ragged2e}
\usepackage{siunitx}
\usepackage{url}
\usepackage[colorlinks=true, urlcolor=blue,  linkcolor=black, citecolor=green]{hyperref}
\usepackage{pdfpages}
\usepackage{blindtext} %Paquete que genera texto ficticio (dummy text) con el comando: \blindtext
\usepackage{cancel} %in the preamble gives you four different modes of striking through: \cancel{text to cancel} draws a diagonal line (slash) through its argument, \bcancel{text to cancel} uses the negative slope (a backslash), \xcancel{text to cancel} draws an X (actually \cancel plus \bcancel), \cancelto{〈value〉}{〈expression〉} draws a diagonal arrow through the 〈expression〉pointing to the 〈value〉 (math-mode only)
\usepackage{csquotes}
\usepackage{afterpage}
\usepackage{parskip} 
\usepackage{float}
\usepackage{enumitem}
\usepackage{multicol}%Paquete que permite la creación aislada de columnas de texto. De esta forma se reduce la cantidad de páginas en nuestro documento
\usepackage{lipsum} %Paquete que genera texto de relleno al igual que \usepackage{blindtext}, se usa con el comando \lipsum[#-#] los #(hashtags) delimitan el número de páginas que deseamos ocupar del paquete lipsum para usarlos como relleno. 

% Definir el comando personalizado para la nota
\newcommand{\mynote}[1]{%
\begin{tikzpicture}[baseline=-0.75ex]
\node[draw, circle, fill=Apple Green!60, inner sep=2pt] (note) {\textbf{Nota:}};
\end{tikzpicture}%
\ \textit{#1}%
}

\newenvironment{Figura}
{\par\medskip\noindent\minipage{\linewidth}}
{\endminipage\par\medskip}
\usepackage{caption}
\usepackage[
backend=biber,
sorting=none,
url=true
]{biblatex} %bibliografía
\addbibresource{biblio.bib}
% Establecer el espacio entre las entradas de la bibliografía
\setlength{\bibitemsep}{1\baselineskip} % Puedes ajustar el valor según tus preferencias
\usepackage{amsmath,amsthm,amssymb,amsfonts}
\usepackage{pifont} %Permite la compilación del comando \xmark: es el símbolo de la tachita
\usepackage{mathtools} %permite la compilación de símbolos para matrices
\usepackage{empheq} %Paquete que se relaciona con \usepackage[most]{tcolorbox} para la creación de cuadros/cajas de colores para delimitar los resultados matemáticos o líneas de texto usando la línea de comando: \begin{empheq}[box={\mymath[colback=orange(webcolor)!70,drop lifted shadow]}]{equation*} \end{empheq}
\usepackage[most]{tcolorbox} %Paquete que se relaciona con \usepackage{empheq} para la creación de cuadros/cajas de colores para delimitar los resultados matemáticos o líneas de texto
\renewcommand{\qedsymbol}{$\blacksquare$}
\usetikzlibrary{positioning,decorations.pathreplacing} %Paquete que permite la compilación de llaves para cuadros sinópticos (brace diagram) 
\usepackage{schemata} %The schemata package is designed just for these brace diagrams
\usepackage{booktabs}

\newcommand\AB[2]{\schema{\schemabox{#1}}{\schemabox{#2}}} %comando o paquete necesario para crear cuadros sinópticos

\newtcbox{\mymath}[1][]{%
nobeforeafter, math upper, tcbox raise base,
enhanced, colframe=blue!30!black,
colback=blue!30, boxrule=1pt,
#1} %Comando importante para encasillar los resultados matemáticos en cajas de diferentes colores, se relaciona con los paquetes \usepackage{empheq} y \usepackage[most]{tcolorbox} 

\newenvironment{sysmatrix}[1]
{\left(\begin{array}{@{}#1@{}}}
{\end{array}\right)}
\newcommand{\ro}[1]{%
\xrightarrow{\mathmakebox[\rowidth]{#1}}%
}
\newlength{\rowidth}% row operation width
\AtBeginDocument{\setlength{\rowidth}{3em}}  %Comando importante para laproducción de lineas de operaciones entre matrices, método de Gauss o Gauss Jordan se relaciona con \newenvironment{sysmatrix}

%formato para cambiar el horario a español
\usepackage[spanish]{datetime2}
\DTMsetdatestyle{spanish}

\renewcommand{\today}{\DTMdisplaydate{\the\year}{\the\month}{\the\day }{-1}}
%%%%%%%%%%%%%%%%%%%%%%%%%%%%%%%%%%%%%%%%%

\usepackage{datetime}
\newcommand{\mycurrenttime}{\xxivtime}

%%%%%%%%%%%%%%%%%%%%%%%%%%%%%%%%%%%%%%%%%%%%%%%%%%%%%%%%%%%%%%%%%%%%%%%%%%%%%%%%%%%%%%%%%%%%%%%%%%%%%%%%%%%%%%%%%%%%%%%%%%%%%%%%%%%%%%%%%%%%%%%%%


%%%%%%%%%%%%%Caja de color para el título%%%%%%%%%%%%%%%%%%%%%%%%%%%%%%%%%%%%%%%%%%%%%%%%%%%%%%

\definecolor{myframecolor}{RGB}{1, 45, 75} % Prussian Blue
\definecolor{myboxcolor}{RGB}{0, 107, 92} % Apple Green

\newtcolorbox{mybox}{
enhanced,
colback=myboxcolor!7, % Color de fondo del cuadro
colframe=myframecolor, % Color del marco del cuadro
arc=0pt,
boxrule=1pt,
borderline west={2mm}{-10mm}{myframecolor}, % Borde en el lado izquierdo
sharp corners=southwest,
width=\linewidth,
before=\par\vspace{\bigskipamount}, % Espacio antes del cuadro
after=\par\vspace{\bigskipamount} % Espacio después del cuadro
}


\newcommand{\euler}{e} %Comando para producir letra e de euler
\newenvironment{remark}{\par\vfill\footnotesize % Vertical white space above the remark and smaller font size
\begin{list}{}{
	\leftmargin=80pt % Indentation on the left
	\rightmargin=60pt}\item\ignorespaces % Indentation on the right
\makebox[-2.5pt]{\begin{tikzpicture}[overlay]
		\node[draw=Horizon!60,line width=2.5pt,circle,fill=Horizon!25,font=\sffamily\bfseries,inner sep=4pt,outer sep=0pt] at (-19pt,5pt){\textcolor{Horizon}{Nota}};\end{tikzpicture}} % Blue Nota in a circle
\advance\baselineskip -1pt}{\end{list}\vskip5pt} % Tighter line spacing and white space after remark
\usepackage{graphicx}

\usepackage{titling}

\input{colores.tex}

%%Producir signos de cita textual``Asignatura''

\newcommand{\xmark}{\ding{55}} %%%comando que genera la tachita

\newenvironment{MyColorPar}[1]{%
\leavevmode\color{#1}\ignorespaces%
}{%
}%

%%%%%%%%%%%%%%%%%%% Título de la tarea, Nombre de alumno

\title{Clase \# 30 - Mecánica Teórica}
\author{\emph{JA} $\boldsymbol{\mid}$}
\date{\today}

%%%%%%%%%%%%%%%%%%% Datos de la Materia

\usepackage{fancyhdr}
\fancypagestyle{plain}{%  the preset of fancyhdr 
\fancyhf{} % clear all header and footer fields
\fancyfoot[R]{
\includegraphics[width=2cm]{LogoFCFMBUAP (1).png}}
\fancyfoot[L]{{\bfseries{\thedate{}}} a las {\bfseries{\mycurrenttime{}}} horas{} (GMT-6, H. Puebla de Zaragoza, Pue)}
\fancyhead[L]{Mecánica teórica}
\fancyhead[R]{\theauthor}
}
\makeatletter
\def\@maketitle{%
\newpage
\null
\vskip 1em%
\begin{center}%
\let \footnote \thanks
{\LARGE \@title \par}%
\vskip 1em%
%{\large \@date}%
\end{center}%
\par
\vskip 1em}
\makeatother

\usepackage{lipsum}  
\usepackage{cmbright}


%%%%%%%%%%%%%%%%%%%%%%%%%%%%%%%%%%%%%%%%%%%%%%%%%%%%%%%%%%%%%%%%%%%%%%%%%%%%%%%%%%%%%%%%%%%%%%%%%%%%%%%%%%%%%%%%%%%%%%%%%%%%%%%%%%%%%%%%%%%%%%%%%%%%%%%%%%%%%%%%%%%%%%%%%%%%%%%%%%%%%%%%%%%%%%%%%%%%%%

\begin{document}



\maketitle


%%%%%%%%%%%%%%%%%%% Datos del profesor, campus y aula

\begin{mybox}
\noindent\begin{tabular}{@{}ll}
	{\bfseries{Profesor}} & Dr. Gilberto Silva Ortigoza\\
	\textcolor{prussianblue}{\bfseries{Campus}} / \textcolor{trueblue}{\bfseries{Facultad}}  & \textcolor{prussianblue}{\bfseries{C.U. BUAP}} /  \textcolor{trueblue}{\bfseries{FCFM}} \\
	\textcolor{prussianblue}{\bfseries{Edificio}} /  \textcolor{trueblue}{\bfseries{EMA4}}    & 	\textcolor{prussianblue}{\bfseries{FM9}} /  \textcolor{trueblue}{\bfseries{302}}   
\end{tabular} 
\end{mybox} \vspace{1cm}

\tableofcontents

\vspace{2cm}

\section{Secciones cónicas}

tenemos:


\[  \frac{1}{r} = \frac{mk}{ \ell^{2} } \left[  1 + \sqrt{  1 + \frac{2E\ell^{2}  }{mk^{2}}  }  \cos (\phi - \bar{\phi}  )\right] \]


Tenemos

\[\frac{1}{r} = C \left[  1 + e \cos (\phi - \bar\phi) \right] \]


\[ C = \frac{mk}{\ell^{2}}; \quad e = \sqrt{ 1 + \frac{ 2E\ell^{2} }{mk^{2}}  }  \]


\begin{enumerate}
\item $e=0 \quad \Rightarrow \quad E = - \frac{mk^{2}}{2\ell^{2}}=E_{0} \quad \textbf{Círculo}$
\item $e<1, \quad E_{0}<E<0 \quad \textbf{Elipse}$
\item $e=1$, $E=0$ \textbf{Parábola}
\item $e>1$, $E>0$ \textbf{Hipérbola}
\end{enumerate}

\begin{figure}[H]
    \centering
	    \includegraphics[width=0.8\textwidth]{1.png} % ajusta el ancho según convenga
    \caption{Descripción clara y concisa de la figura.}
    \label{fig:mi_figura} % opcional: para usar con \ref{fig:mi_figura}
\end{figure}

\textbf{Masa reducida}:

\[  m = \frac{ m_{T} \cdot m_{s} }{ m_{T} + m_{s} } \]

\[ \vec{r} = \vec{r}_{T} - \vec{r}_{S} \]

\textbf{Primera Ley de Kepler:}

\emph{Los planetas se mueven en órbitas elípticas con el sol en uno de sus focos}


Tenemos:


\[  m = \frac{m_{s} m_{T} }{ m_{s}+m_{T} } = \frac{ m_{T} }{1 + \frac{ m_{T} }{m_{s}}}  \]


\begin{equation}
\begin{split}
\left. \begin{array}{cc}
m_{T} &= 5\text{.}972\times 10^{24} kg\\
m_{S} &= 1\text{.}989\times 10^{30} kg
\end{array} \right\} \frac{m_{T} }{m_{s}} \approx 3\text{.}00251\times10^{-6}
\end{split}
\end{equation}

\emph{En este tercer examen si nos aprender de memoria, pero tenemos que sabernos las 3 leyes de Kepler, no deben de venir escritas}.

\textbf{Segunda Ley de Kepler}

\emph{Los planetas barren áreas iguales en tiempos iguales.}

Esta ley está íntimamente ligada al \textbf{momento angular orbital} del planeta. 
En efecto, la constancia del área barrida implica que el momento angular se conserva; 
por tanto, el planeta se mueve en un plano fijo y su velocidad varía de modo que, 
cuando está más cerca del Sol (perihelio), se mueve más rápido, y cuando está más lejos (afelio), más lento.

\smallskip

\textbf{¿Quién determina la forma de la órbita del planeta?}

La forma de la órbita está determinada por la \textbf{energía total del sistema} y el 
\textbf{momento angular orbital}. Dado un valor de momento angular $L$, la energía $E$ 
define el tipo de trayectoria que el cuerpo sigue en el plano:
\[
E < 0 \Rightarrow \text{elipse}, \quad
E = 0 \Rightarrow \text{parábola}, \quad
E > 0 \Rightarrow \text{hipérbola.}
\]
Por tanto, la \textbf{excentricidad $e$} está fijada por los valores de $E$ y $L$.

\smallskip

\textbf{Caso especial: $L = 0$}

Si el momento angular orbital vale cero, eso significa que la posición y la velocidad son 
\textbf{paralelas o antiparalelas}, es decir, el cuerpo se mueve directamente hacia o desde el centro. 
En este caso no existe una órbita en sentido kepleriano: no hay área barrida ni trayectoria elíptica, 
sino un \textbf{movimiento radial puro}.\textbf{Caso especial: $L=0$ — Demostración}

\begin{equation}
\vec{L}=\vec{r}\times m\vec{v}=m r^{2}\,\dot{\theta}\,\hat{\mathbf{z}}
\end{equation}

\begin{equation}
L=0 \quad\Rightarrow\quad r^{2}\dot{\theta}=0
\end{equation}

\begin{equation}
r>0 \ \text{(partícula no en el origen)} \quad\Rightarrow\quad \dot{\theta}=0
\end{equation}

\begin{equation}
\theta(t)=\theta_{0}=\text{cte.} \quad\Rightarrow\quad \vec{r}(t)=r(t)\,\hat{\mathbf{e}}_{r}(\theta_{0})
\end{equation}

\begin{equation}
\vec{v}=\dot{r}\,\hat{\mathbf{e}}_{r}+r\dot{\theta}\,\hat{\mathbf{e}}_{\theta}
\end{equation}

\begin{equation}
\dot{\theta}=0 \quad\Rightarrow\quad \vec{v}=\dot{r}\,\hat{\mathbf{e}}_{r}\ \ (\parallel\ \vec{r})
\end{equation}

\begin{equation}
\Rightarrow\ \ \text{la velocidad es radial (paralela o antiparalela a }\vec{r}\text{): movimiento radial puro.}
\end{equation}

\begin{equation}
\dot{A}=\frac{1}{2}r^{2}\dot{\theta}=\frac{L}{2m}=0 \quad\Rightarrow\quad \text{no se barre área (no hay órbita kepleriana).}
\end{equation}

\begin{equation}
\text{Ecuaciones en polares (fuerza central):}\quad
\begin{cases}
\ddot{r}-r\dot{\theta}^{2}=\dfrac{F(r)}{m}\\[6pt]
r\ddot{\theta}+2\dot{r}\dot{\theta}=0
\end{cases}
\end{equation}

\begin{equation}
\dot{\theta}=0 \ \Rightarrow\ 
\begin{cases}
\ddot{r}=\dfrac{F(r)}{m}\\[6pt]
0=0
\end{cases}
\quad\Rightarrow\quad \text{dinámica 1D puramente radial.}
\end{equation}

\begin{equation}
\text{Energía:}\quad 
E=\frac{1}{2}m\dot{r}^{2}+\underbrace{\frac{L^{2}}{2mr^{2}}}_{\text{término centrífugo}}+V(r)
\end{equation}

\begin{equation}
L=0 \quad\Rightarrow\quad E=\frac{1}{2}m\dot{r}^{2}+V(r)\ \ \text{(sin barrera centrífuga) → caída/escape radial.}
\end{equation}

\begin{equation}
\text{Ley } \frac{1}{r^{2}}:\ \ r(\theta)=\frac{p}{1+e\cos(\theta-\theta_{0})},\quad 
p=\frac{L^{2}}{mk}
\end{equation}

\begin{equation}
L=0 \Rightarrow p=0 \Rightarrow \text{la cónica degenera en la recta }\theta=\theta_{0}\ (\text{trayectoria radial}).
\end{equation}



\begin{figure}[H]
    \centering
	    \includegraphics[width=0.8\textwidth]{2.png} % ajusta el ancho según convenga
    \caption{Descripción clara y concisa de la figura.}
    \label{fig:mi_figura} % opcional: para usar con \ref{fig:mi_figura}
\end{figure}

¿Cómo era la eplise de acuerdo a los datos de nuestra trayectoria (psoiblemente ejercicio en clases)? 

R =  (imagen)


\begin{figure}[H]
    \centering
	    \includegraphics[width=0.8\textwidth]{3.png} % ajusta el ancho según convenga
    \caption{Descripción clara y concisa de la figura.}
    \label{fig:mi_figura} % opcional: para usar con \ref{fig:mi_figura}
\end{figure}


Tenemos:


\textbf{Área infinitesimal}


\begin{figure}[H]
    \centering
	    \includegraphics[width=0.8\textwidth]{4.png} % ajusta el ancho según convenga
    \caption{Descripción clara y concisa de la figura.}
    \label{fig:mi_figura} % opcional: para usar con \ref{fig:mi_figura}
\end{figure}

\textbf{Interpretación del producto cruz}

\[ \left| \vec{A} \times \vec{B}  \right| \]

\begin{figure}[H]
    \centering
	    \includegraphics[width=0.8\textwidth]{5.png} % ajusta el ancho según convenga
    \caption{Descripción clara y concisa de la figura.}
    \label{fig:mi_figura} % opcional: para usar con \ref{fig:mi_figura}
\end{figure}

\begin{figure}[H]
    \centering
	    \includegraphics[width=0.8\textwidth]{6.png} % ajusta el ancho según convenga
    \caption{Descripción clara y concisa de la figura.}
    \label{fig:mi_figura} % opcional: para usar con \ref{fig:mi_figura}
\end{figure}

\textbf{Momento angular orbital}

\begin{equation}
\begin{split}
 \vec{L}&= \vec{r} \times \vec{\mathcal{P}}\\
 &= \vec{r} \times m \frac{ d\vec{r} }{dt}\\
\end{split}
\end{equation}


Tenemos: 


\begin{equation}
\begin{split}
 \vec{L} dt &= m\vec{r} \times d\vec{r}\\
  \frac{\left| \vec{L} dt  \right|}{2} &= \frac{m}{2}  \left|\vec{r} \times d\vec{r} \right|\end{split}
\end{equation}

\textbf{Período orbital}

Cuando barremos toda el área, entonces podremos calcular el período orbital. 

\emph{ Tercera Ley de Kepler }


\[ \tau^{2} \approx a^{3} \]

\textbf{Área de una elipse}: lo que barre la trayectoria del planeta.

\[ A_{elipse} = \pi{ab} \]


\begin{figure}[H]
    \centering
	    \includegraphics[width=0.8\textwidth]{7.png} % ajusta el ancho según convenga
    \caption{Descripción clara y concisa de la figura.}
    \label{fig:mi_figura} % opcional: para usar con \ref{fig:mi_figura}
\end{figure}

\[ \Delta{A} = \frac{ \ell }{2m} \Delta{t}\]


Tenemos:


\[ \boxed{  \tau^{2} = \left( \frac{ 4\pi^{2} m }{k}  \right) a^{3}  } \]



Cuando Newton, se percató de esto pudo percatarse de toda la teoría que desarrolla la evolución temporal de los planetas. 

Esencialmente, hemos terminado con el tema en cuestión.

En este modelo, no existe nada más que la interacción entre el sol y la tierra.

 

\end{document}